\documentclass[10pt,a4paper]{amsart}
\usepackage{amsmath,amssymb,amsthm,graphicx}
\usepackage{verbatim}
\newenvironment{reference}{\comment}{\endcomment}
%\newenvironment{reference}{}{}
\newenvironment{slogan}{\comment}{\endcomment}
\newenvironment{history}{\comment}{\endcomment}

% For commutative diagrams we use Xy-pic
\usepackage[all]{xy}

% We use 2cell for 2-commutative diagrams.
\xyoption{2cell}
% Theorem environments.
%
\theoremstyle{plain}
\newtheorem{theorem}[subsection]{Theorem}
\newtheorem{proposition}[subsection]{Proposition}
\newtheorem{lemma}[subsection]{Lemma}

\theoremstyle{definition}
\newtheorem{definition}[subsection]{Definition}
\newtheorem{example}[subsection]{Example}
\newtheorem{exercise}[subsection]{Exercise}
\newtheorem{situation}[subsection]{Situation}

\theoremstyle{remark}
\newtheorem{remark}[subsection]{Remark}
\newtheorem{remarks}[subsection]{Remarks}

\numberwithin{equation}{subsection}

% Macros
%
\def\lim{\mathop{\mathrm{lim}}\nolimits}
\def\colim{\mathop{\mathrm{colim}}\nolimits}
\def\Spec{\mathop{\mathrm{Spec}}}
\def\Hom{\mathop{\mathrm{Hom}}\nolimits}
\def\Ext{\mathop{\mathrm{Ext}}\nolimits}
\def\SheafHom{\mathop{\mathcal{H}\!\mathit{om}}\nolimits}
\def\SheafExt{\mathop{\mathcal{E}\!\mathit{xt}}\nolimits}
\def\Sch{\mathit{Sch}}
\def\Mor{\mathop{\mathrm{Mor}}\nolimits}
\def\Ob{\mathop{\mathrm{Ob}}\nolimits}
\def\Sh{\mathop{\mathit{Sh}}\nolimits}
\def\NL{\mathop{N\!L}\nolimits}
\def\CH{\mathop{\mathrm{CH}}\nolimits}
\def\proetale{{pro\text{-}\acute{e}tale}}
\def\etale{{\acute{e}tale}}
\def\QCoh{\mathit{QCoh}}
\def\Ker{\mathop{\mathrm{Ker}}}
\def\Im{\mathop{\mathrm{Im}}}
\def\Coker{\mathop{\mathrm{Coker}}}
\def\Coim{\mathop{\mathrm{Coim}}}

% Boxtimes
%
\DeclareMathSymbol{\boxtimes}{\mathbin}{AMSa}{"02}

%
% Macros for moduli stacks/spaces
%
\def\QCohstack{\mathcal{QC}\!\mathit{oh}}
\def\Cohstack{\mathcal{C}\!\mathit{oh}}
\def\Spacesstack{\mathcal{S}\!\mathit{paces}}
\def\Quotfunctor{\mathrm{Quot}}
\def\Hilbfunctor{\mathrm{Hilb}}
\def\Curvesstack{\mathcal{C}\!\mathit{urves}}
\def\Polarizedstack{\mathcal{P}\!\mathit{olarized}}
\def\Complexesstack{\mathcal{C}\!\mathit{omplexes}}
% \Pic is the operator that assigns to X its picard group, usage \Pic(X)
% \Picardstack_{X/B} denotes the Picard stack of X over B
% \Picardfunctor_{X/B} denotes the Picard functor of X over B
\def\Pic{\mathop{\mathrm{Pic}}\nolimits}
\def\Picardstack{\mathcal{P}\!\mathit{ic}}
\def\Picardfunctor{\mathrm{Pic}}
\def\Deformationcategory{\mathcal{D}\!\mathit{ef}}
\usepackage{cleveref}
\crefname{originalreference}{original reference}{original references}
\Crefname{originalreference}{Original reference}{Original references}
\makeatletter
\newlabel{remark-covariance-open-lower-shriek}{{0G56}{source}{}{sourcepointer-1}{}}
\newlabel{remark-covariance-open-lower-shriek@cref}{{[originalreference][0][]0G56}{[1][0][]source}{}{}{}}
\makeatother
\crefname{originallemma}{original lemma}{original lemmas}
\Crefname{originallemma}{Original lemma}{Original lemmas}
\makeatletter
\newlabel{coherent-lemma-Artin-Rees}{{01YA}{source}{}{sourcepointer-2}{}}
\newlabel{coherent-lemma-Artin-Rees@cref}{{[originallemma][0][]01YA}{[1][0][]source}{}{}{}}
\makeatother
\crefname{originallemma}{original lemma}{original lemmas}
\Crefname{originallemma}{Original lemma}{Original lemmas}
\makeatletter
\newlabel{algebra-lemma-intersect-powers-ideal-module-zero}{{00IP}{source}{}{sourcepointer-3}{}}
\newlabel{algebra-lemma-intersect-powers-ideal-module-zero@cref}{{[originallemma][0][]00IP}{[1][0][]source}{}{}{}}
\makeatother
\crefname{originalreference}{original reference}{original references}
\Crefname{originalreference}{Original reference}{Original references}
\makeatletter
\newlabel{algebra-remark-intersection-powers-ideal}{{00IR}{source}{}{sourcepointer-4}{}}
\newlabel{algebra-remark-intersection-powers-ideal@cref}{{[originalreference][0][]00IR}{[1][0][]source}{}{}{}}
\makeatother
\title{Sections with compact support in degree zero are the global sections vanishing near the complement}
\author{The Stacks Project Authors}
\date{stacks-project:duality.tex:local-snapshot-20261012; GNU FDL 1.2 or later}
\begin{document}
\maketitle
\noindent\textit{Editorial note.} All delivered mathematical text is the author's own, copied from the pinned Stacks Project chapter file duality.tex; no statement or proof step was rewritten, reordered or removed.\par\medskip
\noindent\textit{Editorial note.} Preamble provenance: the pinned chapter file opens with an input of the shared preamble and carries no preamble text of its own, so the project's single Stacks-wide preamble staged at sources/stacks-duality/preamble.tex is excerpted verbatim under the recipe fixed by LANE\_BRIEF 22.1, as three spans, lines 8-13 (verbatim package and the reference/slogan/history environments), lines 14-18 (Xy-pic setup) and lines 154-217 (the native theorem declarations and mathematical macros). Nothing was added, removed or redefined; the wrapper class and the amsmath/amssymb/amsthm/graphicx base packages are supplied by the assembler.\par\medskip
\noindent\textit{Editorial note.} The counted statement and proof contain no cross reference, no citation and no included file, so the delivered body is self contained.\par\medskip
\noindent\textit{Editorial note.} Canonical location of the counted proposition: Stacks tag 0G5C (sources/stacks-duality/pinned-original-source.tex line 9276).\par\medskip
\noindent\textit{Allowed standard prerequisites (editorial).}\begin{itemize}
\item The covariance map for the lower shriek and compactly supported cohomology of an open. Same chapter, Remark labelled remark-covariance-open-lower-shriek, cited at proof lines 9294-9296. Provides the map whose image is described.
\item Artin-Rees for the powers of the ideal sheaf. Cohomology of Schemes, Lemma labelled coherent-lemma-Artin-Rees, cited at proof lines 9302-9304. Identifies the limit of the complexes with the limit of the H\textasciicircum{}0(X, I\textasciicircum{}n F).
\item Krull intersection for a finite module over a Noetherian ring. Algebra, Lemma labelled algebra-lemma-intersect-powers-ideal-module-zero and Remark labelled algebra-remark-intersection-powers-ideal, cited at proof lines 9308-9312. Characterises the sections contained in every I\textasciicircum{}n F as those vanishing near Z.
\end{itemize}\medskip
\ifdefined\hypertarget\hypertarget{sourcepointer-1}{}\fi
\noindent\textit{Source reference.} Original reference 0G56: sources/stacks-duality/pinned-original-source.tex, source lines 8923--8923. Same-chapter pointer into the author environment remark at source line 8923; invoked at proof line 9292 of the counted proof. It is retained only as a navigation pointer; every load-bearing step is inside the delivered span.\par\medskip
\ifdefined\hypertarget\hypertarget{sourcepointer-2}{}\fi
\noindent\textit{Source reference.} Original lemma 01YA: sources/stacks-chow/tags.txt, source lines 2549--2549. Cross-chapter standard prerequisite located in the project tag list staged in the local snapshot; invoked at proof line 9301 of the counted proof. It is retained only as a navigation pointer; every load-bearing step is inside the delivered span.\par\medskip
\ifdefined\hypertarget\hypertarget{sourcepointer-3}{}\fi
\noindent\textit{Source reference.} Original lemma 00IP: sources/stacks-chow/tags.txt, source lines 687--687. Cross-chapter standard prerequisite located in the project tag list staged in the local snapshot; invoked at proof line 9308 of the counted proof. It is retained only as a navigation pointer; every load-bearing step is inside the delivered span.\par\medskip
\ifdefined\hypertarget\hypertarget{sourcepointer-4}{}\fi
\noindent\textit{Source reference.} Original reference 00IR: sources/stacks-chow/tags.txt, source lines 689--689. Cross-chapter standard prerequisite located in the project tag list staged in the local snapshot; invoked at proof line 9310 of the counted proof. It is retained only as a navigation pointer; every load-bearing step is inside the delivered span.\par\medskip
\begin{lemma}
\label{lemma-h0-compactly-supported}
Let $X$ be a proper scheme over a field $k$. Let
$K \in D^b_{\textit{Coh}}(\mathcal{O}_X)$ with $H^i(K) = 0$
for $i < 0$. Set $\mathcal{F} = H^0(K)$.
Let $Z \subset X$ be closed with complement $U = X \setminus U$.
Then
$$
H^0_c(U, K|_U) \subset H^0(X, \mathcal{F})
$$
is given by those global sections of $\mathcal{F}$ which
vanish in an open neighbourhood of $Z$.
\end{lemma}

\begin{proof}
Consider the map
$H^0_c(U, K|_U) \to H^0_X(X, K) = H^0(X, K) = H^0(X, \mathcal{F})$ of
Remark \ref{remark-covariance-open-lower-shriek}. To study this
we represent $K$ by a bounded complex $\mathcal{F}^\bullet$ with
$\mathcal{F}^i = 0$ for $i < 0$. Then we have by definition
$$
H^0_c(U, K|_U) = \lim H^0(X, \mathcal{I}^n\mathcal{F}^\bullet)
= \lim \Ker(
H^0(X, \mathcal{I}^n\mathcal{F}^0) \to
H^0(X, \mathcal{I}^n\mathcal{F}^1))
$$
By Artin-Rees (Cohomology of Schemes, Lemma \ref{coherent-lemma-Artin-Rees})
this is the same as $\lim H^0(X, \mathcal{I}^n\mathcal{F})$.
Thus the arrow $H^0_c(U, K|_U) \to H^0(X, \mathcal{F})$
is injective and the image consists of those global sections of $\mathcal{F}$
which are contained in the subsheaf $\mathcal{I}^n\mathcal{F}$
for any $n$. The characterization of these as the sections which
vanish in a neighbourhood of $Z$ comes from Krull's intersection
theorem (Algebra, Lemma \ref{algebra-lemma-intersect-powers-ideal-module-zero})
by looking at stalks of $\mathcal{F}$. See discussion in
Algebra, Remark \ref{algebra-remark-intersection-powers-ideal}
for the case of functions.
\end{proof}

\end{document}
